\begin{exercise}{110}
Show that a monomorphism $i\colon A\to B$ is a weak equivalence if and only if it
is anodyne.

\begin{proof}
First we prove directly that anodyne maps are weak equivalences.
For a horn $\Lambda^n_k\hookrightarrow\Delta^n$, $n\ge1$, use barycentric
coordinates and put
\[
 \mu(t)=\min_{j\ne k}t_j,\qquad
 R(t)_j=
 \begin{cases}
  t_k+n\mu(t),&j=k,\\
  t_j-\mu(t),&j\ne k.
 \end{cases}
\]
The coordinates remain nonnegative with sum $1$, and at least one
coordinate other than $k$ becomes zero.
Thus $R$ retracts the simplex onto the horn, and
$(1-s)t+sR(t)$ is a strong deformation retraction fixing the horn.

The class of monomorphisms whose realizations are homotopy equivalences
is saturated. Indeed, realizations of monomorphisms are CW subcomplex
inclusions (Corollary~1.1.41).
A subcomplex inclusion which is a homotopy equivalence is a strong
deformation retract inclusion, by the homotopy extension property.
Its pushouts remain so, since realization preserves pushouts and the
deformation fixes the attaching subspace.
Coproducts and retracts preserve homotopy equivalences.
For a countable chain, every sphere map, path, and homotopy in the
union has compact image in a finite CW subcomplex, hence lies in
some stage. The map from the first stage therefore induces a
bijection on components and isomorphisms on all homotopy groups;
Whitehead's theorem makes it a homotopy equivalence.
Monomorphisms themselves are saturated by Exercise~\ref{ex:59}.
Since this class contains the horns, Corollary~1.3.16 shows that every
anodyne map is a weak equivalence.

Conversely, suppose $i$ is a monomorphism and a weak equivalence.
Apply Proposition~1.3.9 to the horn inclusions to factor
\[
 A\xrightarrow{j}Z\xrightarrow{p}B,
\]
where $j$ is anodyne and $p$ is a Kan fibration.
The direction just proved makes $j$ a weak equivalence, so $p$ is
also a weak equivalence by 3-for-2.
Exercises~\ref{ex:109} and~\ref{ex:107} make $p$ a trivial fibration.
By Exercise~\ref{ex:59}, it lifts against the monomorphism $i$,
giving $r\colon B\to Z$ with $ri=j$ and $pr=\mathrm{id}_B$.
Consequently
\[
 i\xrightarrow{(\mathrm{id}_A,r)}j
  \xrightarrow{(\mathrm{id}_A,p)}i
\]
is a retract diagram in the arrow category.
Since anodyne maps are closed under retracts, $i$ is anodyne.
\end{proof}
\end{exercise}
